% ============================================================================
% PART 02 - MATHEMATICAL FRAMEWORK
% ============================================================================

\labelsec{math_framework}

This chapter develops the complete mathematical foundation for SHT. We proceed from PDB structure $\rightarrow$ height function $\rightarrow$ local coordinate frame $\rightarrow$ moments $\rightarrow$ geometric units.

\section{From PDB Coordinates to Height Function}

\subsection{PDB Structure Parsing}

\textbf{Input:} PDB file containing atomic coordinates

\textbf{Step 1: Extract DNA backbone}

\begin{align}
    \{P_1, P_2, \ldots, P_n\} &= \text{phosphate atom positions in DNA chain}\\
    r(P_i) &= [x_i, y_i, z_i] \in \mathbb{R}^3 \quad \text{for } i = 1, \ldots, n
\end{align}

Standard DNA residues: DA, DT, DG, DC (or A, T, G, C in RNA contexts)

\subsection{Principal Axis Alignment}

We establish a coordinate system where the DNA helical axis aligns with the z-direction.

\textbf{Step 2: Compute covariance matrix}

\begin{beq}{cov_matrix}
    C = \frac{1}{n} \sum_{i=1}^{n} (\vect{r}_i - \bar{\vect{r}})(\vect{r}_i - \bar{\vect{r}})^T
\end{beq}

where $\bar{\vect{r}} = \frac{1}{n}\sum_i \vect{r}_i$ is the geometric center.

\textbf{Step 3: Eigendecomposition}

\begin{beq}{eigendecomp}
    C = Q \Lambda Q^T
\end{beq}

where
\begin{align}
    \Lambda &= \text{diag}(\lambda_1, \lambda_2, \lambda_3) \quad \text{with } \lambda_1 \geq \lambda_2 \geq \lambda_3 \\
    Q &= [\vect{q}_1, \vect{q}_2, \vect{q}_3] \quad \text{(matrix of eigenvectors)}
\end{align}

\textbf{Step 4: Rotation to align principal axis}

The principal axis (largest variance direction) becomes z-axis:

\begin{beq}{rotation}
    \vect{r}'_i = Q^T \vect{r}_i \quad \text{for all atoms}
\end{beq}

\textbf{Physical meaning:} $\lambda_1 \gg \lambda_2, \lambda_3$ because DNA is elongated along its helical axis. After rotation, the z-coordinates $z'_i$ are the heights.

\subsection{Height Function Smoothing}

Raw phosphate positions are discrete. We smooth them via cubic spline interpolation.

\textbf{Step 5: Arc length parametrization}

\begin{beq}{arclength}
    s_j = \sum_{i=1}^{j-1} \norm{\vect{r}'_i - \vect{r}'_{i-1}} \quad \text{cumulative distance along backbone}
\end{beq}

\textbf{Step 6: Cubic spline}

\begin{align}
    h(s) &= \text{CubicSpline}(s_i, z'_i) \quad \text{with knot spacing } \Delta s = 3.4 \text{ \AA} \\
    &= \text{(one basepair)}
\end{align}

Result: $h(s)$ is smooth, twice-differentiable, and represents height as continuous function of arc length.

\section{Local Coordinate Frames}

\subsection{Frenet-Serret Frames}

Classical differential geometry provides the Frenet-Serret frame along a curve. For height-parametrized curve $\vect{C}(h)$ in 3D:

\begin{beq}{tangent}
    \vect{T}(h) = \frac{d\vect{C}}{dh}
\end{beq}

The tangent vector has components:
\begin{beq}{tangent_components}
    \vect{T}(h) = [\dd{x}{h}, \dd{y}{h}, 1]
\end{beq}

(The z-component is 1 because we parametrize by height $h$.)

\subsection{Height-Parametrized Curvature}

The curvature vector is the second derivative:

\begin{beq}{curvature_vector}
    \vect{\kappa}(h) = \frac{d^2\vect{C}}{dh^2} = [\ddo{x}{h}, \ddo{y}{h}, 0]
\end{beq}

Note: $d^2z/dh^2 = 0$ because we're parametrizing by height.

The magnitude is the curvature scalar:

\begin{beq}{curvature_scalar}
    v_0(h) = \norm{\vect{\kappa}(h)} = \sqrt{(\ddo{x}{h})^2 + (\ddo{y}{h})^2}
\end{beq}

\textbf{Physical interpretation:} This is the inverse radius of instantaneous curvature in the $xy$-plane (perpendicular to z-axis).

\subsection{Comparison with Arc-Length Curvature}

A natural alternative is the classical arc-length (Frenet) curvature:
\begin{beq}{arclength_kappa}
    \kappa(s) = \frac{\norm{\vect{C}'(s) \times \vect{C}''(s)}}{\norm{\vect{C}'(s)}^3}
\end{beq}
which is invariant under rigid-body reparametrization. The two definitions coincide when $ds/dh = 1$, i.e., when the backbone is perfectly aligned with the helical axis. In practice the centroid trajectory $\vect{C}(h)$ departs from this ideal (mean $ds/dh = 0.571 \pm 0.211$ across the N\,=\,200 dataset), so the two quantities need not agree.

To quantify this empirically, we computed both $\bar{\kappa}_h$ and $\bar{\kappa}_s$ from the \emph{same} centroid trajectories $\vect{C}(h)$ for all 200 structures. The results are reported in \S\ref{sec:arclength_comparison}. In brief:
\begin{itemize}
    \item $r(\bar{v}_0, \bar{\kappa}_s) = -0.069$ ($p = 0.33$): the two metrics are statistically independent.
    \item $r(\bar{v}_0, \Omega) = 0.870$ ($p < 10^{-60}$) versus $r(\bar{\kappa}_s, \Omega) = 0.011$ ($p = 0.88$): Vij carries 80$\times$ more structural information than arc-length curvature.
\end{itemize}
These results confirm that the height-parametrization is not a coordinate-choice artifact but captures a qualitatively distinct and biologically more informative aspect of backbone geometry.

\section{Cross-Sectional Coordinate System}

At each height $h$, we establish a 2D coordinate system in the plane perpendicular to $\vec{T}(h)$.

\subsection{Construction of Cross-Section Basis}

\textbf{Step 1: Normalize tangent}

Let $\vec{T}(h)$ be the (possibly non-unit) tangent. Normalize it:

\begin{beq}{tangent_normalized}
    \hat{\vect{T}}(h) = \frac{\vect{T}(h)}{\norm{\vect{T}(h)}}
\end{beq}

\textbf{Step 2: First cross-section axis}

\begin{beq}{cs_axis1}
    \vect{N}_1(h) = \frac{\hat{\vect{T}}(h) \times \vect{z}}{\norm{\hat{\vect{T}}(h) \times \vect{z}}}
\end{beq}

where $\vect{z} = [0, 0, 1]$ is the vertical direction.

This is orthogonal to both the tangent and the z-axis.

\textbf{Step 3: Second cross-section axis}

\begin{beq}{cs_axis2}
    \vect{N}_2(h) = \hat{\vect{T}}(h) \times \vect{N}_1(h)
\end{beq}

By construction, $\{\vect{N}_1, \vect{N}_2\}$ form an orthonormal basis in the cross-sectional plane perpendicular to $\hat{\vect{T}}$.

\subsection{Atom Projection}

For each atom at position $\vect{r}_i$, we project it onto the cross-sectional plane at its corresponding height $h_i$:

\begin{align}
    \vect{r}_{\perp,i} &= \vect{r}_i - (\vect{r}_i \cdot \hat{\vect{T}})\hat{\vect{T}} \quad \text{(remove component along tangent)} \\
    x_i &= \vect{r}_{\perp,i} \cdot \vect{N}_1(h_i) \quad \text{(x-coordinate in cross-section)} \\
    y_i &= \vect{r}_{\perp,i} \cdot \vect{N}_2(h_i) \quad \text{(y-coordinate in cross-section)}
\end{align}

Result: Each atom has $(x, y)$ position in its cross-sectional 2D plane.

\section{Moment Calculations}

\subsection{General Moment Definition}

For a distribution of atoms, the general moment is:

\begin{beq}{general_moment}
    M_{jk} = \sum_{i \in \text{slice}} m_i x_i^j y_i^k
\end{beq}

where $m_i$ is the mass (or unity for unweighted).

\subsection{0th Moment}

\begin{beq}{m0th_moment}
    M_0 = \sum_{i \in \text{slice}} m_i \quad \text{(total mass in slice)}
\end{beq}

\subsection{1st Moments (Center of Mass)}

\begin{align}
    M_{1,0} &= \sum_i m_i x_i \\
    M_{0,1} &= \sum_i m_i y_i
\end{align}

\textit{Note:} After centering the atoms in the cross-section, these should be $\approx 0$.

\subsection{2nd Moment Tensor (Inertia)}

\begin{align}
    I_{xx} &= \sum_i m_i x_i^2 \\
    I_{yy} &= \sum_i m_i y_i^2 \\
    I_{xy} &= \sum_i m_i x_i y_i
\end{align}

The inertia tensor is:

\begin{beq}{inertia_tensor}
    I = \begin{pmatrix} I_{xx} & I_{xy} \\ I_{xy} & I_{yy} \end{pmatrix}
\end{beq}

Eigenvalues $\lambda_1, \lambda_2$ of $I$ are principal moments of inertia.

\subsection{3rd Moment Tensor (Skewness)}

All components:

\begin{align}
    S_{30} &= \sum_i m_i x_i^3 \\
    S_{03} &= \sum_i m_i y_i^3 \\
    S_{21} &= \sum_i m_i x_i^2 y_i \\
    S_{12} &= \sum_i m_i x_i y_i^2
\end{align}

Plus mixed components if needed. The skewness magnitude is:

\begin{beq}{skewness_magnitude}
    \Omega = \sqrt{S_{30}^2 + S_{03}^2 + S_{21}^2 + S_{12}^2}
\end{beq}

\section{Summary of Variables}

\tabref{math_variables} summarizes key variables used throughout.

\begin{table}[h]
    \centering
    \caption{Key Mathematical Variables}\labeltab{math_variables}
    \begin{tabularx}{\textwidth}{lXl}
        \toprule
        \textbf{Symbol} & \textbf{Definition} & \textbf{Units} \\
        \midrule
        $h(s)$ & Height (projection onto principal axis) & \AA \\
        $\vec{C}(h)$ & 3D position as function of height & \AA \\
        $\vec{T}(h)$ & Tangent vector (unnormalized) & dimensionless \\
        $v_0(h)$ & Curvature (Vij) & \AA$^{-1}$ \\
        $\eta_r(h)$ & Flatness (Shambhian) & dimensionless, $\in [0,1]$ \\
        $\Omega(h)$ & Roughness (Mamton) & \AA$^3$ \\
        $M_0$ & 0th moment (total mass) & g/mol \\
        $I$ & Inertia tensor & g$\cdot$\AA$^2$/mol \\
        $S$ & Skewness tensor & g$\cdot$\AA$^3$/mol \\
        \bottomrule
    \end{tabularx}
\end{table}

This completes the mathematical framework. The next chapter derives specific formulas for the three geometric units.
